% TOP_PROBLEM: {"id": 5300012, "title": "Local connectivity of the Mandelbrot set", "category_id": 11, "category": {"id": 11, "name": "dynamical_systems", "display_name": "Dynamical Systems", "description": "Problems about long-term behavior of deterministic systems, Hamiltonian dynamics, and periodic orbits.", "slug": "dynamical-systems", "order_index": 11, "created_at": "2026-07-31T15:26:25.671Z"}, "status": "open"}
% FIRST_POSTED: 2026-09-27
% !TeX program = lualatex
% Compile twice: lualatex 5300012.tex
% All references and prose are embedded; no BibTeX database or external inputs.
\documentclass[11pt,a4paper]{article}
\usepackage[margin=24mm,headheight=14pt]{geometry}
\usepackage{fontspec}
\setmainfont{Latin Modern Roman}
\setsansfont{Latin Modern Sans}
\setmonofont{Latin Modern Mono}
\usepackage{amsmath,amssymb,mathtools}
\usepackage{unicode-math}
\setmathfont{Latin Modern Math}
\usepackage{microtype}
\usepackage{enumitem,xurl,needspace}
\usepackage[dvipsnames]{xcolor}
\usepackage{hyperref}
\hypersetup{unicode=true,colorlinks=true,linkcolor=MidnightBlue,urlcolor=MidnightBlue,
 pdftitle={500 Mathematical Problem Statements},pdfauthor={Source-annotated compilation},bookmarksnumbered=true}
\usepackage{bookmark}
\usepackage{tocloft}
\setlength{\cftsecnumwidth}{3em}
\cftsetpnumwidth{2.5em}
\cftsetrmarg{4em}
\usepackage{titlesec}
\titleformat{\section}{\Large\bfseries\raggedright}{\thesection}{0.7em}{}
\usepackage{fancyhdr}
\pagestyle{fancy}\fancyhf{}
\fancyhead[L]{\small\sffamily Mathematical problem statements}
\fancyhead[R]{\small Source edition 22}
\fancyfoot[C]{\thepage}
\setlength{\parindent}{0pt}
\setlength{\parskip}{5pt plus 1pt}
\setlength{\emergencystretch}{3em}
\setcounter{tocdepth}{1}
\setcounter{secnumdepth}{1}
\setlist[itemize]{leftmargin=3em,itemsep=5pt,topsep=4pt,parsep=0pt}
\allowdisplaybreaks
\newcommand{\N}{\mathbb{N}}
\newcommand{\Z}{\mathbb{Z}}
\newcommand{\Q}{\mathbb{Q}}
\newcommand{\R}{\mathbb{R}}
\newcommand{\C}{\mathbb{C}}
\newcommand{\F}{\mathbb{F}}
\newcommand{\A}{\mathbb{A}}
\newcommand{\PP}{\mathbb P}
\newcommand{\E}{\mathbb{E}}
\newcommand{\eps}{\varepsilon}
\newcommand{\dd}{\,\mathrm d}
\newcommand{\sref}[2]{\hyperlink{source:#1:#2}{[S#2]}}
\newcommand{\eref}[2]{\hyperlink{extra:#1:#2}{[E#2]}}
% Turn the next switch off for a shorter reading copy. The quotations preserve
% every exactTarget field, including qualifications omitted from a short summary.
\newif\ifshowcatalogue
\showcataloguetrue
\newcommand{\cataloguescope}[1]{\ifshowcatalogue
 \paragraph{Exact scope in the supplied catalogue.}
 \begin{quote}\small #1\end{quote}\fi}
\hypersetup{pdftitle={Problem 5300012 - Local connectivity of the Mandelbrot set},pdfauthor={Open Problems Math research notebook}}
\fancyhead[L]{\small\sffamily Open Problems Math \textbar\ Research notebook}
\fancyhead[R]{\small\sffamily 127 \textbar\ 27 September 2026}
\numberwithin{equation}{section}
\begin{document}
\setcounter{section}{0}
% ================================================================
% problemId: 5300012
% Canonical problem ID: 5300012
% exactTarget SHA-256: 0846ad72df6d1c72f53cc923d0076e6aa563fe57147a04e39cff3865b7607470
\clearpage
\section*{\texorpdfstring{Local connectivity of the Mandelbrot set}{Local connectivity of the Mandelbrot set}}\label{5300012}
\subsection{Definitions and mathematical statement}
For $c\in\C$, let $f_c(z)=z^2+c$ and write $f_c^{\circ n}$ for its $n$-fold iterate. The Mandelbrot set is
\[
 \mathcal M=\{c\in\C:\sup_{n\ge0}|f_c^{\circ n}(0)|<\infty\}.
\]
A space is locally connected at a point if every neighborhood contains a connected neighborhood of that point. The question is whether $\mathcal M$ has this property at every $c\in\mathcal M$, particularly on its boundary. Connectedness of the whole set alone is not local connectedness.
\par\smallskip\textit{Formulation and concept references:} \sref{5300012}{1}.
\cataloguescope{Determine whether the Mandelbrot set M = \{c in C : the orbit of 0 under z -> z\textasciicircum{}2 + c is bounded\} is locally connected at every point.}
\subsection{Short English statement}
Is the Mandelbrot set connected at every sufficiently local scale around each of its points?
% BEGIN RESEARCH ADDITION 127
\subsection{Research attempt}

\paragraph{Current frontier.}
Benini's survey distinguishes local connectivity, rigidity, and the
various combinatorial formulations that can be related under additional
hypotheses. \sref{5300012}{1}
Kahn--Kapiamba--Lyubich's June 2026 manuscript establishes a priori
bounds and MLC for a specified class of parabolically bounded primitive
renormalizations; its future-perspective section explicitly describes
remaining unbounded-combinatorics regimes.  It also records progress for
bounded satellite types, so ``satellite'' by itself is not a synonym
for an entirely untreated case. \eref{5300012}{1}, Section 1.
The present route makes the geometric shrinking criterion quantitative
and checks what is lost when dynamical annuli are transferred to
parameter space.

\paragraph{Attack routes.}
One can seek shrinking parameter puzzles, rule out nontrivial fibers,
or transfer dynamical a priori bounds by holonomy.  The first needs
connected neighborhoods in addition to small diameters.  The second
needs a theorem turning the particular fiber statement into local
connectivity, not only uniqueness in a combinatorial class.  The third
needs control of covering degrees and distortion at infinitely many
scales.  We isolate a sufficient annular package with all three
requirements visible, then test whether positivity of individual
moduli can supply it.

\paragraph{Attempt I: a precise shrinking package.}
Fix $c\in\mathcal M$.  Suppose there are bounded Jordan domains
$P_0\supset\overline{P_1}\supset P_1\supset\overline{P_2}\supset\cdots$
with $c\in P_n$ for every $n$.  More simply, require
$\overline{P_{n+1}}\subset P_n$.  Assume also that
\[
 K_n=\overline{P_n}\cap\mathcal M
 \quad\text{is connected for every }n.
\]
Each $K_n$ is a neighborhood of $c$ in the relative topology because
it contains $P_n\cap\mathcal M$.  Put
$A_n=P_n\setminus\overline{P_{n+1}}$ and use
$\operatorname{mod}\{r<|z|<R\}=(2\pi)^{-1}\log(R/r)$.

\textbf{Annular shrinking criterion.} If
\begin{equation}\label{5300012:sum}
 \sum_{n\geq0}\operatorname{mod}(A_n)=\infty,
\end{equation}
then $\mathcal M$ is locally connected at $c$ under the preceding
connected-neighborhood assumption.

Here are the analytic details.  The series inequality for extremal
length of successive essential annuli gives
\[
 \operatorname{mod}(P_0\setminus\overline{P_N})
 \geq\sum_{n=0}^{N-1}\operatorname{mod}(A_n).
\]
One can see the direction directly by putting the extremal metrics on
the disjoint annuli, scaled so that their crossing lengths and areas
both equal their moduli.  Every crossing curve traverses each annulus;
the ratio of squared total length to total area is their sum.
This is the elementary annular form of the Gr\"otzsch inequality;
its role in the recent renormalization argument is discussed in
\eref{5300012}{1}, Section 1.

The nested compact intersection $K=\bigcap_n\overline{P_n}$ is a
continuum.  If it had positive diameter, the doubly connected domain
$P_0\setminus K$ would have finite modulus: a nondegenerate planar
continuum has positive logarithmic capacity, so it is not a puncture
in the conformal classification of rings.  Monotonicity would bound
all the displayed moduli above by this finite value, contradicting
\eqref{5300012:sum}.  Thus $K=\{c\}$.  Nested compactness now gives
$\operatorname{diam}P_n\to0$: otherwise choose two separated points
in arbitrarily late closures and extract limits in $K$, a contradiction.
The connected neighborhoods $K_n$ consequently form a neighborhood
basis at $c$, proving the criterion.

The complex-analysis ring facts in this step are standard extremal-
length input, not new theorems of the notebook; the geometric
hypotheses identifying appropriate \emph{parameter} domains are the
missing dynamical content.

\paragraph{Attempt II: quantify transfer losses.}
Suppose a dynamical annulus $B_n$ of modulus $\mu_n$ is related to an
intermediate annulus $\widetilde B_n$ by an unbranched holomorphic
cover $\widetilde B_n\to B_n$ of degree $d_n$, and $\widetilde B_n$ is mapped to $A_n$ by a
$k_n$-quasiconformal homeomorphism.  Covering an annulus divides its
modulus by the degree, while quasiconformal distortion changes modulus
by a factor between $1/k_n$ and $k_n$.  Hence
\begin{equation}\label{5300012:loss}
 \operatorname{mod}(A_n)\geq\frac{\mu_n}{d_n k_n}.
\end{equation}
For the covering identity, lift both annuli to strips under the
exponential map; the $d_n$-fold winding multiplies the circumference
relative to the strip height.  The distortion bound follows by pulling
back admissible extremal metrics and using the defining quasiconformal
area/length inequality.  These are exactly the normalization-sensitive
steps, not an unexplained transfer of a lower bound.

A concrete sufficient parameter-by-parameter estimate is therefore
\begin{equation}\label{5300012:target}
 \sum_{n\geq0}\frac{\mu_n}{d_n k_n}=\infty,
 \quad\text{together with connectedness of }K_n.
\end{equation}
For instance, $\mu_n\geq\mu_*>0$ and
$d_nk_n\leq C(n+1)$ imply divergence.  The constants may depend on
$c$; the original target asks for local connectivity at each point,
not a uniform modulus of connectivity across all parameters.
Exponential growth of $d_nk_n$, on the other hand, defeats this
particular lower bound even if every dynamical modulus exceeds one
fixed positive number.

The required covering in \eqref{5300012:loss} is explicitly \emph{unbranched
on the annulus}.  A branched return map on a puzzle disk cannot be
inserted without checking where its critical points lie.  Nor is the
parameter holonomy automatically quasiconformal with bounded constants
at all scales.  The a priori bounds and transfer arguments in the
recent manuscript are proved for specified combinatorics; they are
not a universal verification of \eqref{5300012:target}.
\eref{5300012}{1}, Sections 1 and 5 and the appendices.

\paragraph{Attempt III: break weaker versions of the criterion.}
Positivity of every annulus modulus is insufficient.  Take concentric
disks of radii
\[
 r_n=\exp\!\left(-\sum_{j=1}^{n}2^{-j}\right),\qquad r_0=1.
\]
Each successive annulus has modulus $2^{-(n+1)}/(2\pi)>0$, but
$r_n\to e^{-1}$ and the nested intersection is a full disk, not a
point.  Thus saying ``there is an annulus at every renormalization
level'' omits an essential summability issue.

Small diameters also do not by themselves certify local connectivity
if the sets are not connected neighborhoods in the ambient set.
A singleton fiber may be defined using combinatorial separation rather
than a basis of actual connected parameter neighborhoods; its
connection to MLC must be proved in that framework.  The distinction
is emphasized in \sref{5300012}{1}, the discussions of rigidity and fibers.
Our criterion requires connectedness directly so that this step is
not hidden.

Finally, a computer plot cannot establish the annular package at all
levels of a given infinitely renormalizable parameter.  Any finite
collection of positive moduli is compatible with a tail whose sum
converges.  Likewise, failure to construct these particular domains
would not refute local connectivity: \eqref{5300012:target} is sufficient,
not asserted necessary for every possible local-connectivity proof.

\paragraph{Stress test.}
The nested compactness argument is used only after proving that the
intersection is a singleton; compactness alone does not shrink the
disks.  Parameter space and dynamical space are kept distinct, with
the two sources of loss in \eqref{5300012:loss} displayed separately.
The proof does not assume uniformity over all combinatorial types or
all $c$, and it does not deduce connected neighborhoods from a drawing.

The companion computation checks the concentric-annulus example and
the arithmetic of the series lower bound.  The infinite geometric
series and the harmonic-series comparison are proved exactly in the
text; numerical values are not treated as conformal-modulus certificates
for the Mandelbrot set.

\paragraph{Outcome.}
\textbf{Outcome: Conditional reduction.}

The notebook proves an annular shrinking criterion and a quantitative
covering/distortion condition that would imply MLC point by point.
It identifies summability, unbranched covering, and connected parameter
neighborhoods as independent requirements, and gives an explicit
counterexample to replacing summability by positivity.  The full
parameter-domain construction and its estimates for all renormalization
regimes remain unproved.  No new MLC parameter class or priority claim
is asserted.
% END RESEARCH ADDITION 127
\subsection{Sources}
\begin{itemize}\raggedright
\item[\textnormal{[S1]}]\hypertarget{source:5300012:1}{} Anna Miriam Benini, A survey on MLC, Rigidity and related topics, Conformal Geometry and Dynamics 22 (2018).
\url{https://arxiv.org/abs/1709.09869}.
{\small\textit{Catalogue formulation source.}}
% BEGIN RESEARCH REFERENCES 127
\item[\textnormal{[E1]}]\hypertarget{extra:5300012:1}{} Jeremy Kahn, Alex Kapiamba, and Mikhail Lyubich.
\emph{MLC for parabolically bounded primitive renormalization}, arXiv:2606.27272, 25 June 2026. Section 1 (scope, remaining regimes, and the Gr\"otzsch inequality) and Section 5 (geometric and combinatorial transfer estimates). Recent manuscript; its specified combinatorial scope is retained.
\url{https://arxiv.org/html/2606.27272}.
{\small\textit{Consulted 27 September 2026 for the indicated result or scope.}}
% END RESEARCH REFERENCES 127
\end{itemize}


\end{document}
